\documentclass{article}
\usepackage{fontspec}
\setmainfont{Microsoft YaHei}
\XeTeXlinebreaklocale "zh"
\XeTeXlinebreakskip = 0pt plus 1pt
\usepackage[margin=2.2cm]{geometry}
\usepackage{booktabs}
\usepackage{enumitem}
\setlist{nosep}
\title{阶段报告（四） 观测函数病根，朝向误差减半与时变考卷九配方}
\author{}
\date{2026-10-28}
\begin{document}
\maketitle

\section{概要}
本轮先解决观测函数缺陷，再处理时变误差考卷。诊断链七步锁定的病根是深度残差均值的背景离群污染，修复后恒定卷朝向误差从 23.46 度降到 9.19 度，位置误差从 3.28 厘米变到 3.98 厘米。方波考卷九个配方逐个实测，H 配置全程稳定，中位 5.20 厘米与 11.32 度，位置误差比死推算基线低 53\%。overlay 两段视频按新成绩重渲染并补上红色三轴架。指节网格与指节骨架两条路线三次实验全部排除。报告末尾给差距账目，说明离可用红线的距离及其构成。

\section{一 观测函数病根}
诊断从一组反常数字开始。真值位姿下点云对深度图的平均残差是 135 毫米，而位置误差只有 3.28 厘米，两者不相容。逐项排查共七步。系差假设被数据否定，工具点系与手系几乎重合，差 1 度 6 毫米。时差假设被对齐矩阵否定，对角线即最优。标定错配假设被逐帧峰值散布否定，旋转轴散布 21 度，不存在恒定修正。最后一步残差分位数锁定病根，残差中位数约 40 毫米，第 95 分位数达 1 到 2.5 米，少数点云投影采到背景深度，无截尾均值被这批离群值拉爆。修复是把打分函数改成 10\% 截尾均值，去掉最离谱的一成残差再平均。

\section{二 恒定卷结果}
截尾比例做了三点权衡。不截尾时位置 3.28 厘米而朝向 23.46 度。截 20\% 时朝向 10.86 度而位置劣化到 4.80 厘米。截 10\% 时位置 3.98 厘米而朝向 9.19 度，三段误差稳定在 3.6 到 4.9 厘米与 7.7 到 10.9 度之间，无后段发散，定为新最优。朝向误差减掉 61\%，位置只增加 0.7 厘米。渲染层同时补上盲信读数的红色三轴架，三色轴架完整，蓝为真值，红为盲信，绿为滤波。

\section{三 方波考卷九配方}
考卷结构先从数据实测。档位每 25 到 35 帧切换一次，切换点在帧 20，50，85 附近，偏差幅度每关节正负 4.99 度。素材完好性单独验证过，真值投影覆盖率全程 86\% 到 100\%，手距相机 0.6 到 0.8 米，后段失败不是素材问题。

九个配方的成绩见表 \ref{tab:square}。A 恒定假设后段发散。B 时变感知漂移失控，偏差估计冲到 18.9 度。C 与 D 跳变检测重扩散不恢复。A' 与 B' 换成截尾观测后前段改善，B' 中位 5.31 厘米，但中后段仍崩在 15 厘米。F 用评分恶化做跳变信号漏报严重，三次真跳变只抓到一次。G 的档位镜像基于错误假设，偏差漂移远超档位幅度。H 同时改三处，跳变信号改用读数位姿的帧间突变，跳变时读数直推位置瞬间移动 3 厘米以上，这个信号直接且不漏报，偏差估计加高斯先验，标准差 10 度，钳住弱维度漂移。H 的跳变检测精确打在帧 49 与 83，与实测切换点吻合，三段误差 4.5，5.3，4.8 厘米，全程稳定，跳变后 5 到 10 帧恢复。

\begin{table}[h]
\centering
\begin{tabular}{lccccc}
\toprule
配置 & 位置(cm) & 朝向(deg) & 前段(cm) & 中段(cm) & 后段(cm) \\
\midrule
死推算 & 11.08 & 14.38 & 与帧同步 & 与帧同步 & 与帧同步 \\
A 恒定假设 & 5.41 & 12.08 & 3.6 & 3.2 & 11.7 \\
B 时变感知 & 13.94 & 30.46 & 3.8 & 15.0 & 15.0 \\
C 检测重扩散(2deg) & 13.34 & 32.97 & 2.6 & 13.9 & 14.2 \\
D 检测重扩散(5deg) & 13.53 & 35.88 & 与 C 同序 & 与 C 同序 & 与 C 同序 \\
A' 恒定加截尾 & 5.02 & 19.56 & 与 A 同序 & 与 A 同序 & 与 A 同序 \\
B' 时变加截尾 & 5.31 & 12.10 & 3.8 & 15.0 & 15.0 \\
F 截尾加评分检测 & 8.43 & 16.40 & 4.3 & 9.0 & 14.6 \\
G 截尾加档位镜像 & 9.32 & 16.94 & 4.3 & 9.0 & 14.5 \\
H 截尾加读数检测加先验 & 5.20 & 11.32 & 4.5 & 5.3 & 4.8 \\
\bottomrule
\end{tabular}
\caption{方波考卷九配方。H 为最终配置，B' 参数为地板 0.5 度与退火常数 60，加截尾观测，加读数突变检测触发重扩散，加偏差高斯先验。}
\label{tab:square}
\end{table}

机理四条。其一，档位切换周期与滤波收敛时间常数同量级，恒定偏差模型结构上落后半个周期，这是 B'，F，G 中后段崩溃的共同根因。其二，评分恶化作跳变信号不可靠，读数突变才可靠。其三，弱可观测维度的偏差估计会随机游走失控，高斯先验是最小改动的钳制。其四，截尾均值是全套修复的前提，观测判别不干净，后续机制都无从发力。

\section{四 指节路线排除}
朝向改进还试过三条基于手指的路线，全部排除。第一条，假设网格系与实测指节系差一个固定约定旋转，三个指根各自解出的旋转互差 48 到 106 度，单一约定旋转不存在。第二条，把手指网格挂实测链段系，判别裕度在 20 帧真照片上三轴全掉。第三条，绕开网格直接用实测指节位置生成骨架胶囊点云，判别裕度仍掉。原因一致，深度图上手指是细结构且常合拢贴掌，加进点云只会引入噪声。当前点云维持掌块加前臂，这是三次实验后的最优配置。

\section{五 结论与差距账目}
结论三条。其一，观测函数的背景离群污染是几个月来朝向误差卡在 23 度的主因，修复后减到 9.19 度。其二，时变误差可跟踪，正确配置是截尾观测加读数突变检测加偏差先验，方波考卷全程 5.20 厘米。其三，指节路线三次排除，朝向的进一步改进不在这条路上。

差距账目三笔，都有数据。设备账，单目深度相机在 0.6 到 0.8 米距离的固有精度约几毫米到 1 厘米，占去可用红线一大半。模型账，点云模型是掌块加前臂胶囊，没有手指，真值位姿下的平均深度残差中位数约 40 毫米，这是模型近似与深度噪声的物理底。信息账，手腕原地拧转几乎不改变各像素深度，实测判别裕度贴近零，朝向误差里有一部分是照片中不存在的信息。可用红线取位置 1 厘米以内与朝向 2 到 3 度以内，当前成绩差 3 到 4 倍。这个差距不是算法参数能抹平的，需要立体视觉，精确扫描模型与多源融合，属于另一个量级的工作。学术尺度下本轮结果成立，观测病根诊断完整，修复有对照，九个配方有死因分析，全部数值可复算。

\end{document}
